% Generated from locale/tr/content/history/set-theory/mythology.tex; do not hand-edit.


\olfileid[tr]{his}{set}{mythology}

\olsection{Tarihten Çok Mit mi?}

Bu olaylara yüz yılı aşkın bir zamanın ardından geriye bakarken, bu yargıları
sorgulamadan kabul etmemeye dikkat etmeliyiz. Sonuçlar kesinlikle yeni,
heyecan verici ve şaşırtıcıydı. Ama gerçekten ne kadar sarsıcıydılar? Ve
geometrik sezgiye güvenmememiz gerektiğini gerçekten gösteriyorlar mıydı?

Sarsıcılık sorusuyla ilgili olarak \citet{Gouvea2011}, Cantor'un Dedekind'e
yazdığı ünlü ``\emph{je le vois, mais je ne le crois pas}'' notunun
bağlamından epey koparıldığına dikkat çeker. İşte bu bağlamın daha büyük bir
parçası (\citeauthor{Gouvea2011}'den alıntılanmıştır):
\begin{quote}
Nezaketinizden ve sabrınızdan bu kadar çok şey istiyorsam konuya duyduğum
gayreti lütfen bağışlayın; size yakın zamanda gönderdiğim bildirimler benim
için bile öylesine beklenmedik, öylesine yenidir ki sizden, değerli dostum,
doğruluklarına ilişkin bir karar alana dek içim rahat etmeyecek. Benimle aynı
görüşe varmadığınız sürece yalnızca şunu söyleyebilirim: \emph{je le vois,
mais je ne le crois pas.}
\end{quote}
Cantor sonucunun ``öylesine beklenmedik, öylesine yeni'' olduğunu biliyordu.
Ama sonucunu herhangi bir zaman \emph{inanılmaz} bulduğu kuşkuludur.
\citeauthor{Gouvea2011}'in belirttiği gibi, sunduğu ispatı Dedekind'den
denetlemesini istiyordu, o kadar.

Geometrik sezgi sorusuna gelince: Peano uzay dolduran eğrisini hiçbir çizim
eklemeden yayımladı. Fakat Hilbert kendi eğrisini yayımlarken amacını açıkladı:
Okurlar ``aşağıdaki geometrik sezgiden yararlanırlarsa'' onlara Peano'nun
sonucunu açıkça anlayabilecekleri bir yol sunacaktı; ardından tıpkı
\olref[his][set][pathology]{sec}'te verilenler gibi bir dizi \emph{çizim}
ekledi.

Daha genel olarak, yayımlanan ispatlarda çizimler bir ölçüde gözden düşmüş
olsa da matematikçilerin bir şeyleri ispatlarken çizimleri \emph{sık sık}
kullandığı gerçeğinden kaçış yoktur. (Kabaca söylersek iyi matematikçiler
geometrik sezgiye ne zaman güvenebileceklerini bilirler.)

Kısacası: abartıya inanmayın; en azından sorgulamadan kabul etmeyin. Bu konuda
daha fazlası için \citet{Giaquinto2007} kaynağını okuyabilirsiniz.

